\section{知识、符号与组织语义库}

知识这一段看似哲学，其实非常工程。谷雨提到，知识至少可以粗略分成事实性知识、过程性知识和情景性知识。事实性知识回答“是什么”，过程性知识回答“怎么做”，情景性知识记录“发生过什么”。对 Agent 来说，最难的往往不是事实，而是流程和情景：如何申请签证、如何在企业系统里报销、如何按某个团队的惯例修代码、如何根据某个用户的偏好安排旅行。

这就引出 workflow 和 neuro-symbolic。让模型每一步都 plan，灵活但不稳定；把流程写成 code、SOP 或 RPA，稳定但死板。真正可用的系统，很可能要在两者之间来回转化：模型从交互中抽象出 documentation、SOP、workflow、knowledge graph，再在执行时用这些结构约束模型。符号结构不是旧时代遗物，它提供的是一致性、可审计性和组织协同。

\begin{importantbox}{Neuro-symbolic 的新语境}
今天重新谈 neuro-symbolic，不是回到手写规则替代模型，而是让神经模型在开放环境中抽象出可编辑、可检查、可复用的符号结构。模型负责理解模糊现实，符号结构负责稳定流程与共享语义。
\end{importantbox}

随后，概念漂移的讨论把问题从模型带回组织。AGI、Agent、Knowledge、AI-native 这些词会被人和模型反复使用，逐渐失去精确定义。更危险的是，模型会把流行词再生产给人，人再用这些词做决策，形成概念漂移循环。很多团队以为自己在讨论同一个概念，实际上每个人脑子里的对象都不同。

尚晏仪把这个问题说得很产品化：一家公司从零开始，最根本的是语义库。它认为“旅行”是什么，“订机票”是什么，“用户需求”是什么，这些定义会慢慢长出产品、流程和组织语言。如果这些定义不同步，组织协同就会变难；如果 AI 工具不断吸收和生成未经校准的内部文本，长期甚至会破坏组织的理智性。

\begin{warningbox}{热词会偷走判断力}
“AGI”“Agent”“AI-native”“knowledge”这些词如果不被重新定义，就会从思考工具变成思考替代品。团队需要先问：我们在这个场景里到底指什么？成功标准是什么？它对应哪些动作、数据、流程和边界？
\end{warningbox}

谢天宝补充说，symbolic 世界的好处就是定义清楚、一致性强。如果能够连接 symbolic 和 neural，很多一致性问题可以先在 symbolic 层解决，再压回 neural 或 subsymbolic 层。这句话与整场对谈的产品问题其实互相呼应：Agent 不是只要更强模型，也需要更清楚的行动边界、语义定义和组织记忆。

\subsection{本章小结}

知识不只是事实库，还包括流程、情景、偏好和组织语义。Agent 系统如果要长期进入真实组织，就必须处理 neuro-symbolic 的结合：模型负责理解和泛化，符号结构负责一致性和协同。概念漂移不是语言洁癖，而是产品和组织风险；语义库会成为 Agent 公司的一种基础设施。

